\chapter{总结与展望}
\label{cha:conclusion}

本章总结了前四章的研究工作，讨论了本文存在的不足，并展望了下一步工作。

\section{全文总结}

虚拟化环境中，客操作系统的网络 I/O 请求需经由内核网络协议栈逐层封装、多次内存拷贝后才能到达 VirtIO 设备，传统的 socket 路径在协议处理、缓冲区分配和数据拷贝等环节引入了显著的 CPU 开销。针对这一问题，本文以 RISC-V 架构的教学操作系统 rCore 为平台，基于 VirtIO-net 半虚拟化网络设备，设计并实现了一种轻量级的内核旁路网络数据面方案。

本文首先梳理了从传统内核网络栈优化到用户态网络栈（DPDK、AF\_XDP、mTCP 等）的技术演进，指出现有主流内核旁路方案多围绕物理网卡 DMA 模型设计，尚未与半虚拟化环境中的 VirtIO Virtqueue 描述符机制实现原理贯通。在此基础上，本文在内核中分配 17 个物理页帧构建收发环形缓冲区，利用 RISC-V SV39 页表将其同时映射到内核恒等映射地址空间和用户虚拟地址空间，实现了跨特权级的物理内存共享。用户程序在共享 slot 中原地构造完整的以太网帧，内核仅将 slot 的物理地址提交至 Virtqueue 描述符，由 VirtIO 设备通过 DMA 直接读写共享页帧，在收发的内核路径上消除了数据拷贝与动态内存分配。整个方案以约 260 行内核代码和 3 个系统调用实现，未引入任何重量级外部框架。为保证对比的公平性，本文还自行实现了涵盖以太网帧解析、ARP、IPv4/UDP 和 TCP 简易握手的轻量级网络协议处理模块，替换了 rCore 原有的外部依赖，使传统 socket 路径与内核旁路路径共享同一 VirtIO 驱动和物理链路。

实验在 QEMU + SLIRP 环境下，以发送吞吐量和平均往返时延为指标，对内核旁路路径、rCore 传统 socket 路径和 Linux 传统 socket 路径进行了对比测试。结果表明，Bypass 路径相比 rCore Socket 路径发送吞吐量提升了 12\%--35\%，往返时延中位数稳定优于 rCore Socket 路径 7\%--13\%；批量发送分析显示，Bypass 路径的入队阶段仅需 20--104 ticks/pkt，性能瓶颈已从数据拷贝转移到 VirtIO 设备交互本身。

\section{工作展望}

当前方案仍存在若干不足。接收路径采用自旋等待模式，在无数据包到达时持续占用 CPU，后续可引入基于 VirtIO 中断的异步通知机制以提升多任务环境下的 CPU 利用率。此外，当前仅支持单进程独占一组收发缓冲区，多进程场景下需要为每个旁路通道分配独立的共享缓冲区，并利用 VirtIO-net 多队列特性实现多核并行 I/O。安全性方面，共享缓冲区对用户进程完全可读写，后续可将控制头的写入权限分离并在内核侧增加基本的帧格式校验。实验环境方面，所有测试均在 QEMU 用户态模拟环境下完成，VirtIO 设备交互延迟和尾部延迟受模拟开销和 vCPU 调度抖动影响较大，后续可在 KVM 硬件加速或 vhost-net 配置下重新评测，并增加采样次数以获得更稳定的统计结果。最后，当前方案仅验证了 UDP 协议的收发，后续可探索在旁路路径上支持 TCP 协议，通过用户态 TCP 状态机实现可靠传输，进一步拓宽方案的适用范围。

